\chapter{2007年大纲I卷（理）}

\section{单选题}

\begin{problem}
\(\alpha\) 是第四象限角，\(\tan \alpha = -\frac{5}{12}\)，则 \(\sin \alpha =\)（\quad）

\choices
    {\(\frac{1}{5}\)}
    {\(-\frac{1}{5}\)}
    {\(\frac{5}{13}\)}
    {\(-\frac{5}{13}\)}
\end{problem}

\begin{answer}
D
\end{answer}

\begin{solution}
由 \(\tan\alpha=-\frac{5}{12}\)，且 \(\alpha\) 是第四象限角，知 \(\sin\alpha<0\)，\(\cos\alpha>0\). 由 \(1+\tan^2\alpha=\frac{1}{\cos^2\alpha}\) 得 \(\cos\alpha=\frac{12}{13}\)，所以 \(\sin\alpha=\tan\alpha\cos\alpha=-\frac{5}{13}\)，故选 D.
\end{solution}

\begin{problem}
设 \(a\) 是实数，且 \(\frac{a}{1+\i} + \frac{1+\i}{2}\) 是实数，则 \(a =\)（\quad）

\choices
    {\(\frac{1}{2}\)}
    {\(1\)}
    {\(\frac{3}{2}\)}
    {\(2\)}
\end{problem}

\begin{answer}
B
\end{answer}

\begin{solution}
因为 \(a\) 为实数，所以
\[
\frac{a}{1+\i}+\frac{1+\i}{2}=\frac{a(1-\i)}{2}+\frac{1+\i}{2}=\frac{a+1}{2}+\frac{1-a}{2}\i
\]
该数为实数，当且仅当虚部为零，即 \(\frac{1-a}{2}=0\)，所以 \(a=1\)，故选 B.
\end{solution}

\begin{problem}
已知向量 \(\bs{a} = (-5, 6)\)，\(\bs{b} = (6, 5)\)，则 \(\bs{a}\) 与 \(\bs{b}\)（\quad）

\choices
    {垂直}
    {不垂直也不平行}
    {平行且同向}
    {平行且反向}
\end{problem}

\begin{answer}
A
\end{answer}

\begin{solution}
\(\bs{a}\cdot\bs{b}=(-5)\times6+6\times5=0\)，且两个向量都不是零向量，所以 \(\bs{a}\perp\bs{b}\)，故选 A.
\end{solution}

\begin{problem}
已知双曲线的离心率为 2，焦点是 \((-4,0)\)，\((4,0)\)，则双曲线方程为（\quad）

\choices
    {\( \frac{x^{2}}{4}-\frac{y^{2}}{12}=1 \)}
    {\( \frac{x^{2}}{12}-\frac{y^{2}}{4}=1 \)}
    {\( \frac{x^{2}}{10}-\frac{y^{2}}{6}=1 \)}
    {\( \frac{x^{2}}{6}-\frac{y^{2}}{10}=1 \)}
\end{problem}

\begin{answer}
A
\end{answer}

\begin{solution}
双曲线的焦点在 \(x\) 轴上，故设其方程为 \(\frac{x^2}{a^2}-\frac{y^2}{b^2}=1\)，其中 \(c=4\). 由离心率 \(e=\frac{c}{a}=2\)，得 \(a=2\)，再由 \(c^2=a^2+b^2\) 得 \(b^2=16-4=12\). 因此双曲线方程为 \(\frac{x^2}{4}-\frac{y^2}{12}=1\)，故选 A.
\end{solution}

\begin{problem}
设 \(a\)，\(b \in \R\)，集合 \(\{1, a + b, a\} = \left\{0, \frac{b}{a}, b\right\}\)，则 \(b - a =\)（\quad）

\choices
    {\(1\)}
    {\(-1\)}
    {\(2\)}
    {\(-2\)}
\end{problem}

\begin{answer}
C
\end{answer}

\begin{solution}
因为式中出现 \(\frac{b}{a}\)，所以 \(a\ne0\). 右边集合含有 \(0\)，而左边的 \(1\) 不可能为零，故必须有 \(a+b=0\)，即 \(b=-a\). 此时
\(\{1,a+b,a\}=\{1,0,a\}\)，\(\left\{0,\frac ba,b\right\}=\{0,-1,-a\}\).
左边含有 \(1\)，所以右边只能由 \(-a=1\) 取得，从而 \(a=-1\)，\(b=1\). 因此 \(b-a=2\)，故选 C.
\end{solution}

\begin{problem}
下面给出的四个点中，到直线 \(x - y + 1 = 0\) 的距离为 \(\frac{\sqrt{2}}{2}\)，且位于 \(\begin{cases}
x + y - 1 < 0,\\
x - y + 1 > 0
\end{cases}\) 表示的平面区域内的点是（\quad）

\choices
    {\((1,1)\)}
    {\((-1,1)\)}
    {\((-1,-1)\)}
    {\((1,-1)\)}
\end{problem}

\begin{answer}
C
\end{answer}

\begin{solution}
点 \((x,y)\) 到直线 \(x-y+1=0\) 的距离为 \(\frac{|x-y+1|}{\sqrt{2}}\). 由题设距离条件及区域条件 \(x-y+1>0\)，得 \(x-y+1=1\). 对四个点分别检验：\((1,1)\) 满足等式但 \(x+y-1=1>0\)；\((-1,1)\) 有 \(x-y+1=-1\)；\((-1,-1)\) 满足等式且 \(x+y-1=-3<0\)；\((1,-1)\) 有 \(x-y+1=3\). 所以符合条件的点为 \((-1,-1)\)，故选 C.
\end{solution}

\begin{problem}
如图，正四棱柱 \(ABCD-A_{1}B_{1}C_{1}D_{1}\) 中，\(AA_{1} = 2AB\)，则异面直线 \(A_{1}B\) 与 \(AD_{1}\) 所成角的余弦值为（\quad）

\bitmapfigure[max width=0.42\linewidth]{img/2007/syllabus_paper_1_science/q7_fig1.png}

\choices
    {\(\frac{1}{5}\)}
    {\(\frac{2}{5}\)}
    {\(\frac{3}{5}\)}
    {\(\frac{4}{5}\)}
\end{problem}

\begin{answer}
D
\end{answer}

\begin{solution}
设正方形底面边长为 \(s\)，建立直角坐标系，取 \(A=(0,0,0)\)，\(B=(s,0,0)\)，\(D=(0,s,0)\)，\(A_{1}=(0,0,2s)\)，\(D_{1}=(0,s,2s)\). 取两条异面直线的方向向量 \(\overrightarrow{A_{1}B}=(s,0,-2s)\)，\(\overrightarrow{AD_{1}}=(0,s,2s)\)，则
\(\left|\overrightarrow{A_{1}B}\cdot\overrightarrow{AD_{1}}\right|=4s^2\)，\(|\overrightarrow{A_{1}B}|=|\overrightarrow{AD_{1}}|=\sqrt{5}s\).
异面直线所成角取锐角，故其余弦值为 \(\frac{4s^2}{5s^2}=\frac45\)，故选 D.
\end{solution}

\begin{problem}
设 \(a > 1\)，函数 \(f(x) = \log_a x\) 在区间 \([a, 2a]\) 上的最大值与最小值之差为 \(\frac{1}{2}\)，则 \(a =\)（\quad）

\choices
    {\(\sqrt{2}\)}
    {\(2\)}
    {\(2 \sqrt{2}\)}
    {\(4\)}
\end{problem}

\begin{answer}
D
\end{answer}

\begin{solution}
因为 \(a>1\)，所以 \(f(x)=\log_a x\) 在 \([a,2a]\) 上单调递增，故最大值与最小值之差为
\[
\log_a(2a)-\log_a a=\log_a2
\]
由题设 \(\log_a2=\frac12\)，得 \(a^{1/2}=2\)，所以 \(a=4\)，故选 D.
\end{solution}

\begin{problem}
\(f(x)\)，\(g(x)\) 是定义在 \(\R\) 上的函数，\(h(x) = f(x) + g(x)\)，则“\(f(x)\)，\(g(x)\) 均为偶函数”是“\(h(x)\) 为偶函数”的（\quad）

\choices
    {充要条件}
    {充分而不必要的条件}
    {必要而不充分的条件}
    {既不充分也不必要的条件}
\end{problem}

\begin{answer}
B
\end{answer}

\begin{solution}
若 \(f(x)\)，\(g(x)\) 均为偶函数，则对任意 \(x\in\R\)，有 \(f(-x)=f(x)\)，\(g(-x)=g(x)\)，从而 \(h(-x)=f(-x)+g(-x)=f(x)+g(x)=h(x)\)，所以 \(h(x)\) 为偶函数. 反之不成立，例如取 \(f(x)=x\)，\(g(x)=-x\)，则 \(h(x)=0\) 为偶函数，但 \(f(x)\)，\(g(x)\) 都不是偶函数. 因此该条件是充分而不必要的，故选 B.
\end{solution}

\begin{problem}
\(\left(x^{2} - \frac{1}{x}\right)^n\) 的展开式中，常数项为 15，则 \(n =\)（\quad）

\choices
    {\(3\)}
    {\(4\)}
    {\(5\)}
    {\(6\)}
\end{problem}

\begin{answer}
D
\end{answer}

\begin{solution}
展开式的通项中，若从 \(-\frac1x\) 中取 \(k\) 项，则该项为
\[
\mathrm C_n^k(x^2)^{n-k}\left(-\frac1x\right)^k=(-1)^k\mathrm C_n^k x^{2n-3k}
\]
常数项要求 \(2n-3k=0\)，即 \(k=\frac{2n}{3}\) 为整数. 当 \(n=3\) 时 \(k=2\)，常数项系数为 \(\mathrm C_3^2=3\)；当 \(n=4\)，\(5\) 时 \(k\) 不是整数，不存在常数项；当 \(n=6\) 时 \(k=4\)，常数项系数为 \((-1)^4\mathrm C_6^4=15\). 所以 \(n=6\)，故选 D.
\end{solution}

\begin{problem}
抛物线 \(y^{2}=4x\) 的焦点为 \(F\)，准线为 \(l\)，经过 \(F\) 且斜率为 \(\sqrt{3}\) 的直线与抛物线在 \(x\) 轴上方的部分相交于点 \(A\)，\(AK \perp l\)，垂足为 \(K\)，且 \(\triangle AKF\) 的面积是（\quad）

\choices
    {\(4\)}
    {\(3\sqrt{3}\)}
    {\(4\sqrt{3}\)}
    {\(8\)}
\end{problem}

\begin{answer}
C
\end{answer}

\begin{solution}
由 \(y^2=4x\) 知焦点 \(F=(1,0)\)，准线 \(l:x=-1\). 过 \(F\) 的直线方程为 \(y=\sqrt3(x-1)\). 联立抛物线方程得
\(3(x-1)^2=4x\)，\(3x^2-10x+3=0\)
所以交点的横坐标为 \(x=3\) 或 \(x=\frac13\). 位于 \(x\) 轴上方的交点为 \(A=(3,2\sqrt3)\). 因为 \(AK\perp l\)，而 \(l\) 为竖直直线，所以 \(K=(-1,2\sqrt3)\). 于是 \(AK=4\)，点 \(F\) 到直线 \(AK\) 的距离为 \(2\sqrt3\)，故
\[
S_{\triangle AKF}=\frac12\times4\times2\sqrt3=4\sqrt3
\]
故选 C.
\end{solution}

\begin{problem}
函数 \(f(x) = \cos^{2} x - 2\cos^{2} \frac{x}{2}\) 的一个单调增区间是（\quad）

\choices
    {\( \left(\frac{\pi}{3},\frac{2\pi}{3}\right) \)}
    {\( \left(\frac{\pi}{6},\frac{\pi}{2}\right) \)}
    {\(\left(0,\frac{\pi}{3}\right)\)}
    {\(\left(-\frac{\pi}{6},\frac{\pi}{6}\right)\)}
\end{problem}

\begin{answer}
A
\end{answer}

\begin{solution}
由 \(2\cos^2\frac{x}{2}=1+\cos x\)，得 \(f(x)=\cos^2x-1-\cos x\)，所以
\[
f'(x)=-\sin2x+\sin x=\sin x(1-2\cos x)
\]
在 \(\left(\frac\pi3,\frac{2\pi}3\right)\) 内，\(\sin x>0\)，且 \(\cos x<\frac12\)，故 \(f'(x)>0\). 因而该区间是函数的一个单调增区间，故选 A.
\end{solution}

\section{填空题}

\begin{problem}
从班委会 \(5\) 名成员中选出 \(3\) 名，分别担任班级学习委员，文娱委员与体育委员，其中甲，乙二人不能担任文娱委员，则不同的选法共有\fillinblank{} 种.
\end{problem}

\begin{answer}
36
\end{answer}

\begin{solution}
文娱委员只能从其余三人中选出，有 \(3\) 种选法. 文娱委员确定后，学习委员和体育委员从剩余四人中有序选取，选法数为 \(4\times3\). 因此不同的选法共有 \(3\times4\times3=36\) 种.
\end{solution}

\begin{problem}
函数 \(y = f(x)\) 的图象与函数 \(y = \log_3 x\)（\(x > 0\)）的图象关于直线 \(y = x\) 对称，则 \(f(x) =\)\fillinblank{}.
\end{problem}

\begin{answer}
\(3^x\)
\end{answer}

\begin{solution}
关于直线 \(y=x\) 对称的两个函数图象互为反函数. 设 \(y=\log_3x\)，交换 \(x\)，\(y\) 得 \(x=\log_3y\)，从而 \(y=3^x\). 因此 \(f(x)=3^x\).
\end{solution}

\begin{problem}
等比数列 \(\{a_{n}\}\) 的前 \(n\) 项和 \(S_{n}\)，已知 \(S_{1}\)，\(2S_{2}\)，\(3S_{3}\) 成等差数列，则 \(\{a_{n}\}\) 的公比为\fillinblank{}.
\end{problem}

\begin{answer}
\(\frac{1}{3}\)
\end{answer}

\begin{solution}
设等比数列的首项为 \(a_1\)，公比为 \(q\). 则 \(S_1=a_1\)，\(S_2=a_1(1+q)\)，\(S_3=a_1(1+q+q^2)\). 由三项成等差数列，得
\[
S_1+3S_3=4S_2
\]
代入并约去非零的 \(a_1\)，得 \(1+3(1+q+q^2)=4(1+q)\)，即 \(q(3q-1)=0\). 等比数列的公比必须使各项比值有定义，故 \(q\ne0\)，所以 \(q=\frac13\).
\end{solution}

\begin{problem}
一个等腰直角三角形的三个顶点分别在正三棱柱的三条侧棱上. 已知正三棱柱的底面边长为 \(2\)，则该三角形的斜边长为\fillinblank{}.
\end{problem}

\begin{answer}
\(2\sqrt{3}\)
\end{answer}

\begin{solution}
设三棱柱三条侧棱上三个顶点的高度分别为 \(h_1\)，\(h_2\)，\(h_3\). 任意两点在底面上的投影是正三角形的两个顶点，故对应边长的平方分别为 \(4+(h_i-h_j)^2\). 设直角顶点对应的高度为 \(h_1\)，则两条直角边相等给出 \(|h_2-h_1|=|h_3-h_1|\). 若 \(h_2=h_3\)，则边 \(2\) 不可能是两条等边的直角三角形的斜边，故 \(h_2-h_1\) 与 \(h_3-h_1\) 符号相反，记为 \(d\) 和 \(-d\)，于是斜边对应的高度差为 \(2d\). 因而由勾股定理得
\(4+4d^2=2(4+d^2)\)，\(d^2=2\).
所以斜边长的平方为 \(4+4d^2=4+4\times2=12\)，斜边长为 \(2\sqrt3\).
\end{solution}

\section{解答题}

\begin{problem}
设锐角三角形 \(ABC\) 的内角 \(A\)，\(B\)，\(C\) 的对边分别为 \(a\)，\(b\)，\(c\)，\(a = 2b\sin A\).

\begin{enumerate}
\item 求 \(B\) 的大小；
\item 求 \(\cos A + \sin C\) 的取值范围.
\end{enumerate}
\end{problem}

\begin{solution}
由正弦定理，\(\frac{a}{\sin A}=\frac{b}{\sin B}\). 结合 \(a=2b\sin A\)，得 \(\frac{1}{\sin B}=2\)，所以 \(\sin B=\frac12\). 因为三角形 \(ABC\) 为锐角三角形，故 \(B=\frac\pi6\).

又因为 \(A+B+C=\pi\)，所以 \(C=\frac{5\pi}{6}-A\). 锐角条件给出 \(0<A<\frac\pi2\) 且 \(0<C<\frac\pi2\)，从而 \(\frac\pi3<A<\frac\pi2\). 于是
\[
\cos A+\sin C=\cos A+\sin\left(\frac{5\pi}{6}-A\right)=\frac32\cos A+\frac{\sqrt3}{2}\sin A=\sqrt3\cos\left(A-\frac\pi6\right)
\]
当 \(\frac\pi3<A<\frac\pi2\) 时，\(\frac\pi6<A-\frac\pi6<\frac\pi3\)，故
\[
\frac{\sqrt3}{2}<\sqrt3\cos\left(A-\frac\pi6\right)<\frac32
\]
所以 \(\cos A+\sin C\) 的取值范围为 \(\left(\frac{\sqrt3}{2},\frac32\right)\).
\end{solution}

\begin{problem}
某商场经销某商品，根据以往资料统计，顾客采用的付款期为 \(\xi\) 的分布列为

\begin{center}
\small
\begin{tabular}{|c|c|c|c|c|c|}
\hline
\(\xi\) & \(1\) & \(2\) & \(3\) & \(4\) & \(5\) \\
\hline
\(P\) & \(0.4\) & \(0.2\) & \(0.2\) & \(0.1\) & \(0.1\) \\
\hline
\end{tabular}
\end{center}

商场经销一件该商品，采用 \(1\text{ 期}\) 付款，基利润为 \(200\text{ 元}\)；分 \(2\text{ 期}\) 或 \(3\text{ 期}\) 付款，基利润为 \(250\text{ 元}\)；分 \(4\text{ 期}\) 或 \(5\text{ 期}\) 付款，其利润为 \(300\text{ 元}\). \(\eta\) 表示经销一件该商品的利润.

\begin{enumerate}
\item 求事件 \(A\)：“购买该商品的 \(3\) 位顾客中，至少有 \(1\) 位采用 \(1\) 期付款”的概率 \(P(A)\)；
\item 求 \(\eta\) 的分布列及期望 \(E(\eta)\).
\end{enumerate}
\end{problem}

\begin{solution}
假设三位顾客的付款方式相互独立. 每位顾客采用 \(1\) 期付款的概率为 \(0.4\)，因此事件 \(A\) 的对立事件是三位顾客都不采用 \(1\) 期付款，故
\[
P(A)=1-(1-0.4)^3=1-0.6^3=0.784
\]

利润 \(\eta\) 的取值为 \(200\)，\(250\)，\(300\). 其中 \(\eta=200\) 的概率为 \(0.4\)，\(\eta=250\) 的概率为 \(0.2+0.2=0.4\)，\(\eta=300\) 的概率为 \(0.1+0.1=0.2\)，所以其分布列为
\begin{center}
\begin{tabular}{c|ccc}
\(\eta\) & \(200\) & \(250\) & \(300\) \\
\hline
\(P\) & \(0.4\) & \(0.4\) & \(0.2\)
\end{tabular}
\end{center}
从而
\[
E(\eta)=200\times0.4+250\times0.4+300\times0.2=240
\]
\end{solution}

\begin{problem}
四棱锥 \(S-ABCD\) 中，底面 \(ABCD\) 为平行四边形，侧面 \(SBC \perp\) 底面 \(ABCD\). 已知 \(\angle ABC = 45^{\circ}\)，\(AB = 2\)，\(BC = 2\sqrt{2}\)，\(SA = SB = \sqrt{3}\).

\begin{enumerate}
\item 证明：\(SA \perp BC\)；
\item 求直线 \(SD\) 与平面 \(SAB\) 所成角的大小.
\end{enumerate}

\bitmapfigure[max width=0.38\linewidth]{img/2007/syllabus_paper_1_science/q19_fig1.png}
\end{problem}

\begin{solution}
建立直角坐标系，使底面 \(ABCD\) 在 \(xOy\) 平面内，取
\(B=(0,0,0)\)，\(C=(2\sqrt2,0,0)\)，\(A=(\sqrt2,\sqrt2,0)\)，\(D=(3\sqrt2,\sqrt2,0)\).
因为平面 \(SBC\) 垂直于底面，且两平面的交线为 \(BC\)，点 \(S\) 在过 \(BC\) 的垂直平面内，故可设 \(S=(s,0,t)\). 由 \(SA=SB\)，得
\[
(s-\sqrt2)^2+2+t^2=s^2+t^2
\]
所以 \(s=\sqrt2\). 再由 \(SB=\sqrt3\)，得 \(2+t^2=3\)，故 \(t^2=1\). 取 \(t=1\)，则 \(S=(\sqrt2,0,1)\)，取 \(t=-1\) 时以下夹角计算相同.

此时 \(\overrightarrow{SA}=(0,\sqrt2,-1)\)，\(\overrightarrow{BC}=(2\sqrt2,0,0)\)，故 \(\overrightarrow{SA}\cdot\overrightarrow{BC}=0\)，从而 \(SA\perp BC\).

取平面 \(SAB\) 的一个法向量
\[
\bs{n}=\overrightarrow{SA}\times\overrightarrow{SB}=(-\sqrt2,\sqrt2,2)
\]
其中 \(\overrightarrow{SB}=(-\sqrt2,0,-1)\). 又 \(\overrightarrow{SD}=(2\sqrt2,\sqrt2,-1)\)，所以
\(|\bs{n}|=2\sqrt2\)，\(|\overrightarrow{SD}|=\sqrt{11}\)，\(\left|\overrightarrow{SD}\cdot\bs{n}\right|=4\).
设直线 \(SD\) 与平面 \(SAB\) 所成的角为 \(\theta\)，则
\[
\sin\theta=\frac{|\overrightarrow{SD}\cdot\bs{n}|}{|\overrightarrow{SD}|\,|\bs{n}|}=\frac{2}{\sqrt{22}}
\]
因此 \(\theta=\arcsin\frac{2}{\sqrt{22}}\).
\end{solution}

\begin{problem}
设函数 \(f(x) = \e^{x} - \e^{-x}\).

\begin{enumerate}
\item 证明：\(f(x)\) 的导数 \(f'(x) \geqslant 2\)；
\item 若对所有 \(x \geqslant 0\) 都有 \(f(x) \geqslant ax\)，求 \(a\) 的取值范围.
\end{enumerate}
\end{problem}

\begin{solution}
\(f'(x)=\e^x+\e^{-x}\). 由平方恒非负，
\[
\e^x+\e^{-x}-2=\left(\e^{x/2}-\e^{-x/2}\right)^2\geqslant0
\]
故 \(f'(x)\geqslant2\).

对任意 \(x>0\)，题设不等式等价于 \(a\leqslant\frac{f(x)}x\). 令 \(g(x)=f(x)-2x\)，则 \(g'(x)=f'(x)-2\geqslant0\)，且 \(g(0)=0\)，故当 \(x\geqslant0\) 时 \(g(x)\geqslant0\)，即 \(f(x)\geqslant2x\). 所以当 \(a\leqslant2\) 时，\(f(x)\geqslant2x\geqslant ax\)，题设不等式成立. 另一方面，若题设不等式对所有 \(x>0\) 成立，则
\[
a\leqslant\frac{\e^x-\e^{-x}}x
\]
令 \(x\to0^+\)，右端极限为 \(f'(0)=2\)，故必有 \(a\leqslant2\). 因此 \(a\) 的取值范围为 \((-\infty,2]\).
\end{solution}

\begin{problem}
已知椭圆 \(\frac{x^{2}}{3} + \frac{y^{2}}{2} = 1\) 的左，右焦点分别为 \(F_1\)，\(F_2\). 过 \(F_1\) 的直线交椭圆于 \(B\)，\(D\) 两点，过 \(F_2\) 的直线交椭圆于 \(A\)，\(C\) 两点，且 \(AC \perp BD\)，垂足为 \(P\).

\begin{enumerate}
\item 设 \(P\) 点的坐标为 \((x_0, y_0)\)，证明：\(\frac{x_0^{2}}{3} + \frac{y_0^{2}}{2} < 1\)；
\item 求四边形 \(ABCD\) 的面积的最小值.
\end{enumerate}
\end{problem}

\begin{solution}
由椭圆方程知两个焦点为 \(F_1=(-1,0)\)，\(F_2=(1,0)\). 设直线 \(AC\) 的单位方向向量为 \(\bs{u}=(\cos\theta,\sin\theta)\)，则直线 \(BD\) 的单位方向向量可取为 \(\bs{v}=(-\sin\theta,\cos\theta)\). 设 \(P=(x_0,y_0)\)，因为 \(P\) 同时在两条直线上，可写成
\[
P=F_2+t\bs{u}=F_1+s\bs{v}
\]
又因 \(F_2-F_1=(2,0)\)，对上式分别与 \(\bs{u}\) 和 \(\bs{v}\) 做内积，得 \(t=-2\cos\theta\)、\(s=-2\sin\theta\). 从而
\[
P=(1,0)-2\cos\theta(\cos\theta,\sin\theta)
  =(1-2\cos^2\theta,-2\sin\theta\cos\theta)
\]
因此
\[
x_0^2+y_0^2=(1-2\cos^2\theta)^2+4\sin^2\theta\cos^2\theta=1
\]
于是
\[
\frac{x_0^2}{3}+\frac{y_0^2}{2}=1-\frac{2x_0^2}{3}-\frac{y_0^2}{2}<1
\]
从而 \(P\) 在椭圆内部.

设直线 \(AC\) 的单位方向向量为 \((\cos\theta,\sin\theta)\)，则直线 \(BD\) 的单位方向向量可取为 \((-\sin\theta,\cos\theta)\). 先求 \(AC\) 的长度. 过 \(F_2=(1,0)\) 的直线可参数表示为
\[
(x,y)=(1+t\cos\theta,t\sin\theta)
\]
代入椭圆方程，得关于 \(t\) 的方程
\[
\frac{3-\cos^2\theta}{6}t^2+\frac{2\cos\theta}{3}t-\frac23=0
\]
其判别式为 \(\frac43\)，故两根之差的绝对值为 \(\frac{4\sqrt3}{3-\cos^2\theta}\). 因为方向向量是单位向量，所以
\[
AC=\frac{4\sqrt3}{3-\cos^2\theta}
\]
对过 \(F_1\) 且方向为 \((-\sin\theta,\cos\theta)\) 的直线作同样计算，得
\[
BD=\frac{4\sqrt3}{3-\sin^2\theta}=\frac{4\sqrt3}{2+\cos^2\theta}
\]
四边形的两条对角线 \(AC\)，\(BD\) 垂直，因此
\[
S_{ABCD}=\frac12AC\cdot BD=\frac{24}{(3-\cos^2\theta)(2+\cos^2\theta)}
\]
令 \(u=\cos^2\theta\in[0,1]\)，则
\[
(3-u)(2+u)=6+u-u^2=\frac{25}{4}-\left(u-\frac12\right)^2\leqslant\frac{25}{4}
\]
所以 \(S_{ABCD}\geqslant\frac{96}{25}\)，当且仅当 \(u=\frac12\) 时取等号. 四边形 \(ABCD\) 的面积最小值为 \(\frac{96}{25}\).
\end{solution}

\begin{problem}
已知数列 \(\{a_{n}\}\) 中 \(a_{1} = 2\)，\(a_{n+1} = (\sqrt{2} - 1)(a_{n} + 2)\)，\(n = 1\)，\(2\)，\(3\)，\(\cdots\).

\begin{enumerate}
\item 求 \(\{a_{n}\}\) 的通项公式；
\item 若数列 \(\{b_n\}\) 中 \(b_1 = 2\)，\(b_{n+1} = \frac{3b_n + 4}{2b_n + 3}\)，\(n = 1\)，\(2\)，\(3\)，\(\cdots\)，证明：\(\sqrt{2} < b_n \leqslant a_{4n-3}\)，\(n = 1\)，\(2\)，\(3\)，\(\cdots\).
\end{enumerate}
\end{problem}

\begin{solution}
令 \(r=\sqrt2-1\)，则 \(0<r<1\)，且 \(r(\sqrt2+2)=\sqrt2\). 因此
\[
a_{n+1}-\sqrt2=r(a_n+2)-\sqrt2=r(a_n-\sqrt2)
\]
又 \(a_1-\sqrt2=2-\sqrt2=\sqrt2r\)，所以
\(a_n-\sqrt2=\sqrt2r^n\)，\(a_n=\sqrt2(1+r^n)\).

记 \(\varphi(x)=\frac{3x+4}{2x+3}\). 先证下界. 直接计算得
\[
\varphi(x)-\sqrt2=\frac{(3-2\sqrt2)(x-\sqrt2)}{2x+3}=\frac{r^2(x-\sqrt2)}{2x+3}
\]
由于 \(b_1=2>\sqrt2\)，归纳可得 \(b_n>\sqrt2\) 对所有正整数 \(n\) 成立.

又 \(\varphi'(x)=\frac{1}{(2x+3)^2}>0\)，所以 \(\varphi\) 在 \((\sqrt2,+\infty)\) 上单调递增. 对任意正整数 \(k\)，由 \(a_k=\sqrt2(1+r^k)\) 得
\[
\varphi(a_k)-\sqrt2=\frac{\sqrt2r^{k+2}}{3+2\sqrt2(1+r^k)}
\]
因为 \(3+2\sqrt2=r^{-2}\)，故
\(3+2\sqrt2(1+r^k)\geqslant r^{-2}\)，\(\varphi(a_k)-\sqrt2\leqslant\sqrt2r^{k+4}=a_{k+4}-\sqrt2\).
即 \(\varphi(a_k)\leqslant a_{k+4}\).

当 \(n=1\) 时，\(b_1=2=a_1\). 若 \(b_n\leqslant a_{4n-3}\)，则由 \(\varphi\) 的单调性及上面的结论，\(b_{n+1}=\varphi(b_n)\leqslant\varphi(a_{4n-3})\leqslant a_{4n+1}=a_{4(n+1)-3}\). 归纳得 \(b_n\leqslant a_{4n-3}\). 结合下界，结论为 \(\sqrt2<b_n\leqslant a_{4n-3}\).
\end{solution}
